\chapter{Algebry tensorowa, symetryczna i zewnętrzna}
\label{ch:lin11-algebry}
\index{algebra tensorowa@Algebra tensorowa}
\index{algebra symetryczna@Algebra symetryczna}
\index{algebra zewnętrzna@Algebra zewnętrzna}
\index{iloczyn zewnętrzny@Iloczyn zewnętrzny}

\noindent\textbf{Cel rozdziału.}
Iloczyn tensorowy z Rozdziału~\ref{ch:lin10-iloczyny-tensorowe}
umożliwia linearyzację odwzorowań wieloliniowych.
W tym rozdziale łączymy wszystkie potęgi tensorowe
w jedną algebrę. Dwa ilorazy tej algebry pozwalają
wymusić przemienność czynników albo znikanie iloczynu
dwóch jednakowych wektorów. W ten sposób powstaną
algebry symetryczna i zewnętrzna. Wyprowadzimy ich
własności uniwersalne, bazy i wymiary, a także
związek algebry zewnętrznej z wyznacznikiem
i minorami macierzy.

Przez cały rozdział \(K\) oznacza ciało, \(0_K\ne1_K\).
Nie zakładamy, że \(\operatorname{char}K=0\).
Własności konstrukcyjne dotyczą dowolnych przestrzeni
wektorowych nad \(K\); przy wyliczaniu wymiarów
wyraźnie zakładamy skończony wymiar. Przyjmujemy
\(V^{\otimes0}=K\) oraz \(V^{\otimes1}=V\), a w potęgach
wyższych ustalamy lewostronne nawiasowanie.
Kanoniczne izomorfizmy łączności z
Twierdzenia~\ref{thm:lin10-lacznosc} umożliwiają
przechodzenie między różnymi nawiasowaniami.

\section{Algebry stopniowane i potęgi tensorowe}
\label{sec:lin11-podstawy}

\begin{definicja}[Algebra, ideał i stopniowanie]
\label{def:lin11-algebra}
\emph{Łączną algebrą z jedynką nad \(K\)}
nazywamy przestrzeń wektorową \(A\) nad \(K\)
z dwuliniowym mnożeniem
\(\mu\colon A\times A\to A\), dla którego
\((ab)c=a(bc)\), oraz z elementem \(1_A\)
spełniającym \(1_Aa=a1_A=a\).
Nie wymagamy przemienności mnożenia.
\emph{Homomorfizm algebr z jedynką}
jest odwzorowaniem liniowym zachowującym
mnożenie i element jednostkowy.

Podprzestrzeń \(I\subseteq A\)
nazywamy \emph{ideałem obustronnym},
jeśli \(aI\subseteq I\) i \(Ia\subseteq I\)
dla każdego \(a\in A\).
Algebra \(A\) jest \emph{nieujemnie stopniowana},
jeśli
\(A=\bigoplus_{k\geq0} A^k\) jako przestrzeń
wektorowa i \(A^pA^q\subseteq A^{p+q}\).
W sumie prostej dopuszczamy tylko elementy
o skończenie wielu niezerowych składowych.
\end{definicja}

\begin{stwierdzenie}[Algebra ilorazowa]
\label{prop:lin11-iloraz-algebry}
Niech \(A\) będzie łączną algebrą z jedynką,
a \(I\subseteq A\) ideałem obustronnym.
Na przestrzeni \(A/I\) istnieje dokładnie jedno
mnożenie spełniające
\[
 (a+I)(b+I)=ab+I\qquad(a,b\in A).
\]
Jest ono dwuliniowe i łączne; jedynką jest
\(1_A+I\). Jeśli \(I=A\), iloraz jest
algebrą zerową, której jedynka jest zerem.
Jeżeli \(A\) jest stopniowana, a
\(I=\bigoplus_{k\geq0}(I\cap A^k)\), to
\[
 A/I\cong\bigoplus_{k\geq0}
       A^k/(I\cap A^k)
\]
jako algebra stopniowana.
\end{stwierdzenie}
\begin{proof}
Jeśli \(a'=a+i\), \(b'=b+j\)
dla \(i,j\in I\), to
\(a'b'-ab=aj+ib+ij\in I\).
Mnożenie klas jest więc niezależne od
reprezentantów. Dwuliniowość i łączność
przechodzą z \(A\) do ilorazu;
\(1_A+I\) spełnia obie tożsamości jednostki.
Jednoznaczność wynika ze wzoru na iloczyn klas.

W przypadku stopniowanym każdą klasę
reprezentuje skończona suma
jednorodnych składników. Jeśli
\(\sum_k a_k\in I\), jednorodność ideału
wymusza \(a_k\in I\cap A^k\)
dla każdego \(k\). To dowodzi zarówno
rozpiętości, jak i jednoznaczności
rozkładu klas na składowe stopni.
Mnożenie respektuje stopnie,
ponieważ \(A^pA^q\subseteq A^{p+q}\).
\end{proof}

\begin{definicja}[Algebra tensorowa]
\label{def:lin11-tensorowa}
Dla przestrzeni wektorowej \(V\) nad \(K\)
\emph{algebrą tensorową} nazywamy
stopniowaną przestrzeń
\[
 T(V):=\bigoplus_{k=0}^{\infty}V^{\otimes k}.
\]
Mnożenie na składowych jednorodnych
\(x\in V^{\otimes p}\),
\(y\in V^{\otimes q}\) definiujemy
przez kanoniczne odwzorowanie
\[
 V^{\otimes p}\otimes_K V^{\otimes q}
 \longrightarrow V^{\otimes(p+q)},
 \qquad x\otimes y\longmapsto xy,
\]
odpowiadające złączeniu czynników.
W szczególności
\[
 (v_1\otimes\cdots\otimes v_p)
 (w_1\otimes\cdots\otimes w_q)
 =v_1\otimes\cdots\otimes v_p
  \otimes w_1\otimes\cdots\otimes w_q.
\]
Mnożenie przedłużamy dwuliniowo na
skończone sumy składowych.
Elementem jednostkowym jest \(1_K\in V^{\otimes0}\).
\end{definicja}

\begin{twierdzenie}[Własność uniwersalna algebry tensorowej]
\label{thm:lin11-tensorowa-uniwersalna}
Dla dowolnej przestrzeni \(V\)
algebra \(T(V)\) jest łączną algebrą z jedynką.
Dla każdej łącznej algebry \(A\) z jedynką
i każdego odwzorowania liniowego
\(f\colon V\to A\)
istnieje dokładnie jeden homomorfizm
algebr z jedynką
\(\widehat f\colon T(V)\to A\) taki, że
\(\widehat f(v)=f(v)\) dla każdego \(v\in V\).
Na tensorach prostych ma on postać
\[
 \widehat f(v_1\otimes\cdots\otimes v_k)
 =f(v_1)\cdots f(v_k),\qquad
 \widehat f(a)=a1_A\quad(a\in V^{\otimes0}=K).
\]
\end{twierdzenie}
\begin{proof}
Mnożenie jednorodnych tensorów jest dobrze
określone dzięki własności uniwersalnej
iloczynu tensorowego i kanonicznym
izomorfizmom łączności
(Twierdzenie~\ref{thm:lin10-lacznosc}).
Na tensorach prostych oba nawiasowania
\((xy)z\) i \(x(yz)\) zapisują czynniki
w identycznej kolejności.
Ponieważ tensory proste rozpinają
każdą potęgę tensorową, równość rozszerza
się na wszystkie składowe jednorodne,
a następnie przez dwuliniowość na \(T(V)\).
Analogicznie konkatenacja z tensorem
pustym daje \(1_Kx=x1_K=x\).
Skończoność liczby niezerowych składowych
gwarantuje, że każdy otrzymany iloczyn
należy do algebraicznej sumy prostej.

Ustalmy \(f\colon V\to A\).
Dla \(k\geq1\) odwzorowanie
\((v_1,\ldots,v_k)\mapsto f(v_1)\cdots f(v_k)\)
jest \(k\)-liniowe: przy ustaleniu
pozostałych argumentów jest złożeniem
odwzorowania liniowego \(f\) i
mnożenia przez ustalone elementy \(A\).
Twierdzenie~\ref{thm:lin10-wielokrotny}
daje jedyne liniowe
\(f_k\colon V^{\otimes k}\to A\)
o wskazanych wartościach.
W stopniu \(0\) określamy
\(f_0(a)=a1_A\).
Dla \(t=\sum_{k=0}^{N}t_k\)
definiujemy
\(\widehat f(t)=\sum_{k=0}^{N}f_k(t_k)\).
Rozkład na składowe jest jednoznaczny,
więc definicja jest dobra i liniowa.
Na dowolnych tensorach prostych mamy
\(\widehat f(xy)=\widehat f(x)\widehat f(y)\);
dwuliniowość przenosi tę równość
na wszystkie pary \(x,y\in T(V)\).
Ponadto \(\widehat f(1_K)=1_A\),
więc \(\widehat f\) jest homomorfizmem
z jedynką.

Jeśli \(g\colon T(V)\to A\) jest drugim
takim homomorfizmem,
to \(g(v_1\otimes\cdots\otimes v_k)
=g(v_1)\cdots g(v_k)
=f(v_1)\cdots f(v_k)\).
W stopniu \(0\) każdy homomorfizm
z jedynką spełnia \(g(a)=a1_A\).
Ponieważ tensory proste i składowa
\(K\) rozpinają \(T(V)\),
otrzymujemy \(g=\widehat f\).
\end{proof}

\begin{przyklad}[Dlaczego algebra tensorowa nie jest przemienna]
\label{prz:lin11-nieprzemienna}
Niech \(V=\mathbb R^2\) z bazą
\((e_1,e_2)\), a \(A=\operatorname{Mat}_{2\times2}(\mathbb R)\).
Zadajmy \(f(e_1)=E_{12}\),
\(f(e_2)=E_{21}\), gdzie \(E_{ij}\)
są jednostkami macierzowymi.
Własność uniwersalna daje homomorfizm
\(\widehat f\colon T(V)\to A\), zatem
\[
 \widehat f(e_1e_2)=E_{11},\qquad
 \widehat f(e_2e_1)=E_{22}.
\]
Wyniki są różne, więc
\(e_1e_2\ne e_2e_1\) w \(T(V)\).
Porządek czynników zachowuje informację.
Przemienność można dopiero wymusić,
przechodząc do odpowiedniego ilorazu.
\end{przyklad}

\section{Algebra symetryczna}
\label{sec:lin11-symetryczna}

\begin{definicja}[Algebra symetryczna]
\label{def:lin11-symetryczna}
Niech \(J\) będzie ideałem obustronnym
\(T(V)\) generowanym przez wszystkie
\(u\otimes v-v\otimes u\), gdzie \(u,v\in V\).
\emph{Algebrą symetryczną} nazywamy
\[
 S(V):=T(V)/J.
\]
Iloczyn klas oznaczamy zwykłym mnożeniem.
Składową stopnia \(k\) oznaczamy
\(S^k(V):=V^{\otimes k}/(J\cap V^{\otimes k})\).
W szczególności \(S^0(V)=K\),
\(S^1(V)=V\), a \(S(V)\) jest
algebrą przemienną.
\end{definicja}

\begin{stwierdzenie}[Stopniowanie i przemienność]
\label{prop:lin11-sym-stopnie}
Ideał \(J\) jest jednorodny, więc
\(S(V)=\bigoplus_{k\geq0}S^k(V)\).
Dla wszystkich \(u,v\in V\) w ilorazie
zachodzi \(uv=vu\); mnożenie w \(S(V)\)
jest przemienne, łączne i posiada jedynkę.
\end{stwierdzenie}
\begin{proof}
Każdy generator \(u\otimes v-v\otimes u\)
leży w \(V^{\otimes2}\).
Dowolny element ideału \(J\)
jest skończoną sumą elementów
\(a(u\otimes v-v\otimes u)b\).
Rozkładając \(a,b\) na jednorodne
składowe, otrzymujemy składniki
jednorodne należące do \(J\).
Stąd \(J=\bigoplus_{k\ge0}(J\cap V^{\otimes k})\).
Twierdzenie~\ref{prop:lin11-iloraz-algebry}
daje stopniowanie i łączność.
W ilorazie obrazy generatorów
\(u,v\in V\) komutują.
Każdy tensor prosty jest iloczynem
takich generatorów; zamiany sąsiednich
czynników pozwalają przestawić
dowolne dwa jednomiany. Z liniowości
wynika przemienność wszystkich
elementów \(S(V)\).
Ponieważ \(J\) ma składniki
wyłącznie stopnia co najmniej \(2\),
klasy w stopniach \(0,1\) nie są
dodatkowo utożsamiane.
\end{proof}

\begin{twierdzenie}[Własność uniwersalna algebry symetrycznej]
\label{thm:lin11-sym-uniwersalna}
Niech \(A\) będzie przemienną,
łączną algebrą z jedynką nad \(K\)
i \(f\colon V\to A\) odwzorowaniem
liniowym. Istnieje dokładnie jeden
homomorfizm algebr z jedynką
\(\overline f\colon S(V)\to A\)
taki, że \(\overline f(v)=f(v)\)
dla wszystkich \(v\in S^1(V)=V\).
\end{twierdzenie}
\begin{proof}
Z Twierdzenia~\ref{thm:lin11-tensorowa-uniwersalna}
otrzymujemy jedyny homomorfizm
\(\widehat f\colon T(V)\to A\)
przedłużający \(f\).
Ponieważ mnożenie w \(A\)
jest przemienne,
\[
 \widehat f(u\otimes v-v\otimes u)
 =f(u)f(v)-f(v)f(u)=0.
\]
Jądro homomorfizmu jest
ideałem obustronnym, a więc
zawiera ideał \(J\) generowany
przez te różnice.
Wobec tego wzór
\(\overline f(t+J)=\widehat f(t)\)
jest dobrze określony, liniowy
i zachowuje mnożenie oraz jedynkę.
Przedłuża \(f\).
Jeśli \(g\) jest innym takim
homomorfizmem, to
\(g(v_1\cdots v_k)
=f(v_1)\cdots f(v_k)\)
i \(g(a)=a1_A\) dla \(a\in K\).
Ponieważ takie iloczyny
rozpinają \(S(V)\), mamy
\(g=\overline f\).
\end{proof}

\begin{twierdzenie}[Baza potęg symetrycznych]
\label{thm:lin11-sym-baza}
Niech \(V\) będzie skończenie
wymiarową przestrzenią nad \(K\),
\(n=\dim V\), z uporządkowaną
bazą \(B=(e_1,\ldots,e_n)\).
Dla każdego \(k\ge0\)
rodzina jednomianów
\[
 e_1^{a_1}\cdots e_n^{a_n},
 \qquad a_i\in\mathbb N_0,\quad
 \sum_{i=1}^n a_i=k,
\]
jest bazą przestrzeni \(S^k(V)\).
Gdy \(n\ge1\), zachodzi
\[
 \dim S^k(V)=\binom{n+k-1}{k}.
\]
Gdy \(n=0\), przestrzeń \(S^0(V)=K\),
a \(S^k(V)=\{0\}\) dla \(k\ge1\).
Wszystkie te twierdzenia obowiązują
w dowolnej charakterystyce ciała.
\end{twierdzenie}
\begin{proof}
Każdy tensor prosty jest iloczynem
wektorów \(V\), a każdy z nich
jest kombinacją liniową \(e_i\).
Po rozwinięciu przez wieloliniowość
i przestawieniu czynników w \(S(V)\)
otrzymujemy kombinację wskazanych
jednomianów; rozpinają one \(S^k(V)\).

Dla niezależności zbudujmy
homomorfizm
\(P\colon S(V)\to K[x_1,\ldots,x_n]\).
Niech \(f(e_i)=x_i\);
jest to dobrze określone
odwzorowanie liniowe \(V\to K[x_1,\ldots,x_n]\).
Własność uniwersalna
z Twierdzenia~\ref{thm:lin11-sym-uniwersalna}
daje homomorfizm \(P\)
spełniający
\(P(e_1^{a_1}\cdots e_n^{a_n})
=x_1^{a_1}\cdots x_n^{a_n}\).
Z definicji pierścienia
wielomianów różne jednomiany
są liniowo niezależne nad \(K\).
Zatem jeśli kombinacja wskazanych
wektorów w \(S^k(V)\) jest zerowa,
to wszystkie jej współczynniki
są zerowe po zastosowaniu \(P\).
To dowodzi niezależności.

Dla \(n\ge1\) liczba
\(n\)-tek nieujemnych liczb całkowitych
o sumie \(k\) jest
\(\binom{n+k-1}{k}\):
można je zakodować,
umieszczając \(n-1\) przegród
wśród \(n+k-1\) miejsc.
Nie używaliśmy dzielenia
przez żaden element ciała,
więc twierdzenie obejmuje
również dodatnią charakterystykę.
Przy \(n=0\) jedyna krotka
wykładników jest pusta
i ma sumę \(0\); daje
bazę \(S^0(V)=K\).
Dla \(k>0\) nie ma
żadnego jednomianu,
co zgodnie z rozpiętością
oznacza przestrzeń zerową.
\end{proof}

\begin{wniosek}[Model wielomianowy]
\label{cor:lin11-wielomiany}
Przy założeniach
Twierdzenia~\ref{thm:lin11-sym-baza}
istnieje izomorfizm algebr
\[
 S(V)\cong K[x_1,\ldots,x_n],
 \qquad e_i\longmapsto x_i,
\]
zachowujący stopnie.
Dla \(n=0\) prawa strona
oznacza \(K\).
\end{wniosek}
\begin{proof}
Homomorfizm \(P\) skonstruowany
w poprzednim dowodzie
przekształca bazę każdej
składowej stopnia \(k\)
bijektywnie na bazę jednomianów
jednorodnych stopnia \(k\).
Jest więc bijekcją liniową
na każdej składowej, a zatem
na algebraicznej sumie prostej.
Jako homomorfizm z jedynką
jest izomorfizmem algebr.
\end{proof}

\begin{przyklad}[Potęga symetryczna przestrzeni dwuwymiarowej]
\label{prz:lin11-sym2}
Niech \(V=K^2\) z bazą \((e_1,e_2)\).
Wówczas
\(S^2(V)=\operatorname{span}
 \{e_1^2,e_1e_2,e_2^2\}\)
i \(\dim S^2(V)=3\).
W \(S(V)\) mamy \(e_1e_2=e_2e_1\),
podczas gdy w \(V\otimes V\)
oba tensory są różne.
W charakterystyce \(2\)
zachodzi \((e_1+e_2)^2=e_1^2+e_2^2\);
nie oznacza to zaniku
elementu \(e_1e_2\), który
nadal należy do wskazanej bazy.
\end{przyklad}


\section{Algebra zewnętrzna i jej własność uniwersalna}
\label{sec:lin11-zewnetrzna}

\begin{definicja}[Algebra zewnętrzna]
\label{def:lin11-zewnetrzna}
Niech \(I\) będzie ideałem
obustronnym \(T(V)\) generowanym przez
wszystkie elementy \(v\otimes v\),
gdzie \(v\in V\).
\emph{Algebrą zewnętrzną} przestrzeni \(V\)
nazywamy iloraz
\[
 \Lambda(V):=T(V)/I.
\]
Iloczyn klas oznaczamy przez \(\wedge\);
obraz tensora prostego
\(v_1\otimes\cdots\otimes v_k\)
zapisujemy jako
\(v_1\wedge\cdots\wedge v_k\).
Składową stopnia \(k\) oznaczamy przez
\[
 \Lambda^k(V):=
 V^{\otimes k}/(I\cap V^{\otimes k}).
\]
W szczególności
\(\Lambda^0(V)=K\) oraz \(\Lambda^1(V)=V\).
\end{definicja}

\begin{stwierdzenie}[Stopniowanie i reguła znaków]
\label{prop:lin11-wedge-znaki}
Algebra \(\Lambda(V)\) jest łączną
algebrą stopniowaną z jedynką.
Dla wszystkich \(u,v\in V\)
zachodzą równości
\[
 v\wedge v=0,\qquad
 u\wedge v=-v\wedge u.
\]
Jeżeli \(\alpha\in\Lambda^p(V)\),
\(\beta\in\Lambda^q(V)\), gdzie
\(p,q\geq0\), to
\[
 \alpha\wedge\beta
 =(-1)^{pq}\beta\wedge\alpha.
\]
W charakterystyce \(2\)
druga równość dla wektorów
przyjmuje postać
\(u\wedge v=v\wedge u\), ale
warunek \(v\wedge v=0\)
nadal obowiązuje.
\end{stwierdzenie}
\begin{proof}
Generatory ideału \(I\) mają
stopień \(2\), a pomnożenie
ich przez elementy jednorodne
z obu stron daje elementy
o określonym stopniu.
Po rozłożeniu ogólnych
mnożników na skończone sumy
jednorodne otrzymujemy
\(I=\bigoplus_k(I\cap V^{\otimes k})\).
Stąd
Twierdzenie~\ref{prop:lin11-iloraz-algebry}
zapewnia stopniowanie,
łączność i jednostkę.
Generatory ideału mają stopień
co najmniej \(2\), więc
\(I\cap K=I\cap V=\{0\}\).

Z definicji \(v\wedge v=0\).
Po zastosowaniu tej równości
do \(u+v\) i rozwinięciu
za pomocą dwuliniowości mamy
\[
 0=(u+v)\wedge(u+v)
  =u\wedge v+v\wedge u.
\]
Daje to regułę zmiany znaku
bez dzielenia przez \(2\).
Na elementach prostych
\(\alpha=v_1\wedge\cdots\wedge v_p\)
i \(\beta=w_1\wedge\cdots\wedge w_q\)
przejście od \(\alpha\wedge\beta\)
do \(\beta\wedge\alpha\) wymaga
\(pq\) zamian kolejnych
czynników sąsiednich.
Każda zamiana wnosi minus,
stąd czynnik \((-1)^{pq}\).
Wielowektory proste rozpinają
składowe jednorodne; dwuliniowość
rozszerza wzór na dowolne
\(\alpha,\beta\).
Jeśli \(p=0\) lub \(q=0\),
zachodzi zwykła przemienność
mnożenia przez skalar.
\end{proof}

\begin{uwaga}[Dlaczego warunek kwadratu jest istotny]
\label{rem:lin11-char2}
Nad ciałem o charakterystyce
różnej od \(2\) z tożsamości
\(v\wedge w=-w\wedge v\)
dla \(v=w\) wynika
\(2(v\wedge v)=0\), a więc
\(v\wedge v=0\).
W charakterystyce \(2\)
ten wniosek jest niedozwolony.
Dlatego algebra zewnętrzna
jest zdefiniowana przez ideał
generowany \emph{kwadratami}
\(v\otimes v\), a nie tylko
przez elementy
\(u\otimes v+v\otimes u\).
W ten sposób naprzemienność
zachodzi nad każdym ciałem.
\end{uwaga}

\begin{twierdzenie}[Własność uniwersalna algebry zewnętrznej]
\label{thm:lin11-zew-uniwersalna}
Niech \(A\) będzie łączną algebrą
z jedynką nad \(K\), a
\(f\colon V\to A\) odwzorowaniem
liniowym spełniającym
\(f(v)^2=0\) dla każdego \(v\in V\).
Wtedy istnieje dokładnie jeden
homomorfizm algebr z jedynką
\(\overline f\colon\Lambda(V)\to A\)
o \(\overline f(v)=f(v)\)
dla wszystkich \(v\in\Lambda^1(V)=V\).
\end{twierdzenie}
\begin{proof}
Własność uniwersalna \(T(V)\)
daje dokładnie jeden homomorfizm
\(\widehat f\colon T(V)\to A\)
przedłużający \(f\).
Dla generatora ideału \(I\)
otrzymujemy
\(\widehat f(v\otimes v)=f(v)^2=0\).
Ponieważ jądro \(\widehat f\)
jest ideałem obustronnym,
zawiera cały \(I\).
Możemy więc określić
\(\overline f(t+I)=\widehat f(t)\).
Jeśli \(t+I=t'+I\), to
\(t-t'\in I\subseteq\ker\widehat f\),
więc obrazy są równe.
Odwzorowanie jest zatem
dobrze określone; zachowuje
liniowość, mnożenie i jedynkę.

Jeśli \(g\) jest drugim
homomorfizmem z wymaganym
ograniczeniem do \(V\), to
\[
 g(v_1\wedge\cdots\wedge v_k)
 =f(v_1)\cdots f(v_k).
\]
W stopniu \(0\) musi zachodzić
\(g(a)=a1_A\). Ponieważ
takie iloczyny wraz ze skalarami
rozpinają \(\Lambda(V)\),
otrzymujemy \(g=\overline f\).
\end{proof}

\begin{twierdzenie}[Własność uniwersalna potęgi zewnętrznej]
\label{thm:lin11-potega-wedge-uniwersalna}
Niech \(k\geq1\).
Odwzorowanie
\[
 \tau_k\colon V^k\to\Lambda^k(V),\qquad
 \tau_k(v_1,\ldots,v_k)
 =v_1\wedge\cdots\wedge v_k
\]
jest \(k\)-liniowe i naprzemienne.
Dla każdej przestrzeni \(Z\)
oraz każdego naprzemiennego
\(k\)-liniowego odwzorowania
\(b\colon V^k\to Z\)
istnieje dokładnie jedno
odwzorowanie liniowe
\(\overline b\colon\Lambda^k(V)\to Z\)
takie, że
\[
 \overline b(v_1\wedge\cdots\wedge v_k)
 =b(v_1,\ldots,v_k).
\]
Teza pozostaje prawdziwa
dla dowolnej charakterystyki \(K\).
\end{twierdzenie}
\begin{proof}
Wieloliniowość \(\tau_k\)
wynika z wieloliniowości
kanonicznego iloczynu tensorowego
i liniowości ilorazu.
Jeśli dwa argumenty są równe,
przestawiamy czynniki aż
znajdą się obok siebie.
Każda zamiana wprowadza minus,
a sąsiednie równe czynniki
mają iloczyn zewnętrzny zero.
Otrzymujemy więc \(\tau_k=0\)
przy powtórzeniu argumentu.

Niech \(b\) będzie naprzemienne.
Z
Twierdzenia~\ref{thm:lin10-wielokrotny}
istnieje dokładnie jedno
liniowe
\(L_b\colon V^{\otimes k}\to Z\)
z
\(L_b(v_1\otimes\cdots\otimes v_k)
=b(v_1,\ldots,v_k)\).
Udowodnimy, że
\(I\cap V^{\otimes k}\subseteq\ker L_b\).
Ideał \(I\) jest rozpięty
przez elementy
\(a(v\otimes v)c\), gdzie
\(a,c\) są tensorami dowolnych,
lecz skończonych stopni.
Po wybraniu jednorodnej
składowej stopnia \(k\)
i rozwinięciu \(a,c\) na
tensory proste otrzymujemy
elementy postaci
\[
 x\otimes v\otimes v\otimes y,
 \qquad
 x\in V^{\otimes p},\
 y\in V^{\otimes q},\
 p+q=k-2.
\]
Dla \(p=0\) lub \(q=0\)
odpowiedni czynnik jest skalarem.
Dla każdego takiego elementu
wartość \(L_b\) jest równa
wartości \(b\) na krotce,
która ma dwa sąsiednie,
identyczne argumenty \(v\);
wynosi więc zero.
Stąd faktycznie
\(I\cap V^{\otimes k}\subseteq\ker L_b\).

Na mocy własności uniwersalnej
ilorazu
(Twierdzenie~\ref{thm:lin2-wlasnosc-ilorazu})
\(L_b\) przechodzi na
\(\overline b\colon\Lambda^k(V)\to Z\)
o wymaganej wartości.
Wielowektory proste rozpinają
\(\Lambda^k(V)\), gdyż są
obrazami tensorów prostych
rozpinających \(V^{\otimes k}\).
Dlatego ich wartości
wyznaczają \(\overline b\)
jednoznacznie.
W dowodzie nie dzieliliśmy
przez żaden skalar.
\end{proof}

\begin{figure}[!htbp]
\centering
\begin{tikzpicture}[>=Latex,node distance=28mm and 37mm]
 \node (prod) {\(V^k\)};
 \node[right=of prod] (wedge) {\(\Lambda^k(V)\)};
 \node[below=of wedge] (z) {\(Z\)};
 \draw[->] (prod)--node[above]{\(\tau_k\)} (wedge);
 \draw[->] (prod)--node[below left]{\(b\)} (z);
 \draw[->] (wedge)--node[right]{\(\overline b\)} (z);
\end{tikzpicture}
\caption{Linearyzacja naprzemiennego odwzorowania
\(k\)-liniowego. Zewnętrzna potęga
spełnia własność uniwersalną niezależnie
od charakterystyki ciała.}
\label{fig:lin11-zewnetrzna-uniwersalna}
\end{figure}

\section{Baza i współrzędne potęg zewnętrznych}
\label{sec:lin11-baza-wedge}

\begin{twierdzenie}[Baza potęg zewnętrznych]
\label{thm:lin11-zew-baza}
Niech \(\dim V=n<\infty\)
i \(B=(e_1,\ldots,e_n)\)
będzie uporządkowaną bazą \(V\).
Dla każdego \(k\ge1\)
rodzina
\[
 e_{i_1}\wedge\cdots\wedge e_{i_k},
 \qquad
 1\leq i_1<\cdots<i_k\leq n,
\]
stanowi bazę przestrzeni \(\Lambda^k(V)\).
Zatem dla \(0\leq k\leq n\)
mamy
\[
 \dim\Lambda^k(V)=\binom nk,
\]
a dla \(k>n\) mamy
\(\Lambda^k(V)=\{0\}\).
Przy \(k=0\) bazą jest
jednoelementowa rodzina \((1_K)\).
W szczególności
\[
 \dim\Lambda(V)=2^n.
\]
\end{twierdzenie}
\begin{proof}
Każdy \(v\in V\) ma jednoznaczny
zapis \(v=\sum_{i=1}^n x_ie_i\).
Po rozwinięciu iloczynu
\(v_1\wedge\cdots\wedge v_k\)
przez wieloliniowość otrzymujemy
sumę elementów
\(e_{j_1}\wedge\cdots\wedge e_{j_k}\).
Jeśli jakiś indeks się powtarza,
iloczyn znika, bo powtórzony
wektor można przenieść na
sąsiednie miejsce.
Pozostałe indeksy możemy
uporządkować rosnąco, zmieniając
znak przy każdej zamianie.
Ponieważ wielowektory proste
rozpinają \(\Lambda^k(V)\),
wskazana rodzina również ją rozpina.

Aby udowodnić niezależność,
nie wystarczy samo porządkowanie.
Dla każdej rosnącej krotki
\(I=(i_1<\cdots<i_k)\)
określmy formę
\[
 D_I(v_1,\ldots,v_k)
 :=\det\bigl(e^{i_a}(v_b)\bigr)_{a,b=1}^k,
\]
gdzie \(e^1,\ldots,e^n\)
jest bazą dualną.
Wyznacznik jest liniowy
w każdej kolumnie, a znika,
gdy dwie kolumny są równe.
Zatem \(D_I\) jest naprzemienną
formą \(k\)-liniową.
Twierdzenie~\ref{thm:lin11-potega-wedge-uniwersalna}
daje jedyny funkcjonał
\(\overline D_I\colon\Lambda^k(V)\to K\)
z
\(\overline D_I(v_1\wedge\cdots\wedge v_k)
=D_I(v_1,\ldots,v_k)\).
Na wektorach naszej rodziny
bazowej mamy
\[
 \overline D_I(e_{j_1}\wedge\cdots\wedge e_{j_k})
 =\det(\delta_{i_a j_b})_{a,b=1}^k
 =\begin{cases}
   1_K,&(i_1,\ldots,i_k)=(j_1,\ldots,j_k),\\
   0_K,&\text{w przeciwnym przypadku}.
  \end{cases}
\]
Jeśli kombinacja wskazanych
wielowektorów jest zerowa,
zastosowanie kolejnych
\(\overline D_I\) daje
zerowanie każdego współczynnika.
Rodzina jest więc liniowo
niezależna, a zatem jest bazą.

Rosnących krotek \(k\) indeksów
ze zbioru \(n\)-elementowego
jest \(\binom nk\).
Gdy \(k>n\), żadna taka
krotka nie istnieje, a
udowodniona rozpiętość oznacza,
że cała przestrzeń jest zerowa.
Dla \(k=0\) składowa jest równa
\(K\), więc ma wymiar \(1\).
Dodając wymiary składowych,
otrzymujemy
\(\dim\Lambda(V)=\sum_{k=0}^n\binom nk
=2^n\). Przypadek \(n=0\)
również jest objęty argumentem:
\(\Lambda(V)=K\).
\end{proof}

\begin{wniosek}[Współczynniki iloczynu zewnętrznego]
\label{cor:lin11-wedge-minory}
Przy założeniach
Twierdzenia~\ref{thm:lin11-zew-baza}
niech \(1\leq k\leq n\)
oraz
\(v_b=\sum_{i=1}^n a_{ib}e_i\)
dla \(b=1,\ldots,k\).
Wówczas
\[
 v_1\wedge\cdots\wedge v_k
 =\sum_{1\leq i_1<\cdots<i_k\leq n}
   \det(a_{i_a b})_{a,b=1}^k\,
   e_{i_1}\wedge\cdots\wedge e_{i_k}.
\]
\end{wniosek}
\begin{proof}
Z
Twierdzenia~\ref{thm:lin11-zew-baza}
wiemy, że rozpatrywane
wielowektory bazowe są niezależne.
Zastosowanie funkcjonału
\(\overline D_I\), skonstruowanego
w tamtym dowodzie, do obu stron
rozwinięcia wyodrębnia współczynnik
przy \(e_{i_1}\wedge\cdots\wedge e_{i_k}\).
Jest on równy
\(D_I(v_1,\ldots,v_k)
=\det(a_{i_a b})_{a,b=1}^k\).
Ponieważ dzieje się tak dla
każdego \(I\), wskazany wzór
jest jedynym możliwym
rozwinięciem współrzędnych.
\end{proof}

\begin{przyklad}[Dwuwektor i jego minory]
\label{prz:lin11-dwuwektor}
W \(\mathbb R^3\) z bazą
standardową niech
\(u=e_1+2e_2\) oraz
\(v=e_2+e_3\).
Wtedy
\[
 u\wedge v
 =e_1\wedge e_2+
  e_1\wedge e_3+
  2e_2\wedge e_3.
\]
Macierz kolumn \([u\ v]\)
ma postać
\(\begin{pmatrix}1&0\\2&1\\0&1\end{pmatrix}\).
Jej minory wierszowe
dla indeksów \((1,2)\),
\((1,3)\), \((2,3)\)
wynoszą odpowiednio
\(1,1,2\). Są to dokładnie
współczynniki otrzymanego
dwuwektora.
\end{przyklad}

\begin{przyklad}[Nie każdy dwuwektor jest prosty]
\label{prz:lin11-nieprosty}
Niech \(V=\mathbb R^4\)
z bazą standardową i
\(\omega=e_1\wedge e_2+
e_3\wedge e_4\in\Lambda^2(V)\).
Z reguły znaków otrzymujemy
\[
 \omega\wedge\omega
 =2e_1\wedge e_2\wedge e_3\wedge e_4
 \ne0,
\]
bo kwadraty składników prostych
znikają, a składniki mieszane
są równe. Natomiast dla
każdego \(u,v\in V\)
\((u\wedge v)\wedge(u\wedge v)=0\),
ponieważ \(u\) i \(v\)
występują dwukrotnie.
Zatem \(\omega\) nie jest
pojedynczym iloczynem
zewnętrznym \(u\wedge v\).
Argument przez niezerowy
kwadrat jest podany nad
\(\mathbb R\), gdzie \(2\ne0\);
nie przenosimy go w tej
postaci do charakterystyki \(2\).
\end{przyklad}


\section{Potęgi zewnętrzne odwzorowań i wyznacznik}
\label{sec:lin11-operator-wedge}

\begin{twierdzenie}[Potęga zewnętrzna odwzorowania]
\label{thm:lin11-potega-operatora}
Niech \(V,W\) będą przestrzeniami
wektorowymi nad \(K\), a
\(A\colon V\to W\) odwzorowaniem liniowym.
Dla każdego \(k\ge1\) istnieje dokładnie
jedno odwzorowanie liniowe
\[
 \Lambda^k A\colon\Lambda^k(V)\to\Lambda^k(W)
\]
spełniające
\[
 (\Lambda^k A)(v_1\wedge\cdots\wedge v_k)
 =Av_1\wedge\cdots\wedge Av_k.
\]
Dla \(k=0\) przyjmujemy
\(\Lambda^0 A=\operatorname{id}_K\).
Jeśli \(B\colon W\to U\) jest liniowe,
to dla każdego \(k\ge0\)
\[
 \Lambda^k(B\circ A)
 =(\Lambda^k B)\circ(\Lambda^k A),
 \qquad
 \Lambda^k(\operatorname{id}_V)
 =\operatorname{id}_{\Lambda^k(V)}.
\]
\end{twierdzenie}
\begin{proof}
Dla \(k\ge1\) funkcja
\((v_1,\ldots,v_k)\mapsto
 Av_1\wedge\cdots\wedge Av_k\)
jest \(k\)-liniowa, ponieważ
\(A\) jest liniowe, a iloczyn
zewnętrzny jest wieloliniowy.
Jeśli dwa argumenty są równe,
ich obrazy przez \(A\) również
są równe, więc iloczyn znika.
Funkcja jest naprzemienna,
a
Twierdzenie~\ref{thm:lin11-potega-wedge-uniwersalna}
zapewnia istnienie i
jednoznaczność \(\Lambda^k A\).

Dla \(k\ge1\) obie strony
pierwszej równości o złożeniu
wysyłają
\(v_1\wedge\cdots\wedge v_k\)
na \(BA v_1\wedge\cdots\wedge BA v_k\).
Te wielowektory rozpinają
\(\Lambda^k(V)\), zatem
równość obowiązuje na całej
przestrzeni. Dowód własności
dla identyczności jest analogiczny.
Przy \(k=0\) oba wzory są
równościami odwzorowań
identycznościowych na \(K\).
\end{proof}

\begin{twierdzenie}[Minory jako współczynniki macierzy]
\label{thm:lin11-minory}
Niech \(n,m\in\mathbb N_0\),
\(A=(a_{ij})\in\operatorname{Mat}_{m\times n}(K)\)
będzie macierzą odwzorowania
\(A\colon K^n\to K^m\).
Niech \(k\ge1\).
W rosnąco uporządkowanych bazach
potęg zewnętrznych współczynnik
macierzy \(\Lambda^k A\)
w wierszu
\(I=(i_1<\cdots<i_k)\)
i kolumnie
\(J=(j_1<\cdots<j_k)\)
jest równy
\[
 [\Lambda^k A]_{I,J}
 =\det(a_{i_a j_b})_{a,b=1}^{k}.
\]
Jeśli \(k>m\) lub \(k>n\),
odpowiednia baza jest pusta,
a macierz ma zerową liczbę
wierszy lub kolumn.
\end{twierdzenie}
\begin{proof}
Dla wektora bazowego
\(e_J=e_{j_1}\wedge\cdots\wedge e_{j_k}\)
mamy
\[
 (\Lambda^k A)e_J
 =Ae_{j_1}\wedge\cdots\wedge Ae_{j_k}.
\]
Zapisując
\(Ae_{j_b}=\sum_{i=1}^m a_{i j_b}e_i\)
i korzystając z
Wniosku~\ref{cor:lin11-wedge-minory},
odczytujemy współczynnik przy
\(e_{i_1}\wedge\cdots\wedge e_{i_k}\):
jest dokładnie wyznacznikiem
macierzy \((a_{i_a j_b})\).
Jest on więc elementem
macierzy operatora w zadanych
bazach. Dla pustych baz
nie ma odpowiednich
wierszy lub kolumn do obliczenia.
\end{proof}

\begin{wniosek}[Wzór Cauchy'ego--Bineta]
\label{cor:lin11-cauchy-binet}
Niech \(m,n,p\in\mathbb N_0\),
\(A\in\operatorname{Mat}_{m\times n}(K)\),
\(B\in\operatorname{Mat}_{n\times p}(K)\).
Niech \(k\) będzie całkowite
i \(0\le k\le\min(m,p)\).
Dla rosnących krotek indeksów
\(I=(i_1<\cdots<i_k)\)
z \(\{1,\ldots,m\}\) oraz
\(J=(j_1<\cdots<j_k)\)
z \(\{1,\ldots,p\}\)
zachodzi
\[
 \det((AB)_{I,J})
 =\sum_{\substack{H\subseteq\{1,\ldots,n\}\\|H|=k}}
     \det(A_{I,H})\det(B_{H,J}).
\]
Dla \(k=0\) obie strony
wynoszą \(1_K\).
Jeśli \(k>n\), suma po
pustej rodzinie indeksów
wynosi \(0_K\), zgodnie z tezą.
\end{wniosek}
\begin{proof}
Dla \(k\ge1\) z
Twierdzenia~\ref{thm:lin11-potega-operatora}
wynika
\(\Lambda^k(AB)=(\Lambda^k A)(\Lambda^k B)\).
W
Twierdzeniu~\ref{thm:lin11-minory}
wykazaliśmy, że
elementami macierzy
tych odwzorowań są odpowiednie
minory rozmiaru \(k\).
Z definicji mnożenia
macierzy element w
wierszu \(I\) i kolumnie
\(J\) jest sumą po wszystkich
rosnących \(k\)-tkach \(H\)
indeksów przestrzeni pośredniej.
Porównując elementy obu
stron równości operatorowej,
otrzymujemy żądaną tożsamość.
Jeśli \(k>n\), przestrzeń
\(\Lambda^k(K^n)\) jest zerowa,
więc odwzorowanie złożone
również jest zerowe; prawa
strona ma postać sumy pustej.
Dla \(k=0\) stosujemy konwencję
wyznacznika macierzy pustej
i potęgi \(\Lambda^0 A=\operatorname{id}_K\).
Obie strony są wtedy równe \(1_K\).
\end{proof}

\begin{przyklad}[Sprawdzenie wzoru Cauchy'ego--Bineta]
\label{prz:lin11-cauchy}
Niech
\[
 A=\begin{pmatrix}1&2&0\\0&1&1\end{pmatrix},
 \qquad
 B=\begin{pmatrix}1&0\\0&1\\1&2\end{pmatrix}.
\]
Dla \(k=2\) jedyne
rosnące indeksy przestrzeni
pośredniej \(\mathbb R^3\)
to \(12,13,23\).
Odpowiednie minory wynoszą
\[
 \begin{aligned}
 (\det A_{I,H})_H&=(1,1,2),\\
 (\det B_{H,J})_H&=(1,2,-1).
 \end{aligned}
\]
Zatem
\(\det(AB)=1\cdot1+1\cdot2+
2\cdot(-1)=1\).
Bezpośrednio
\(AB=\begin{pmatrix}1&2\\1&3\end{pmatrix}\),
której wyznacznik również
wynosi \(1\).
\end{przyklad}

\begin{twierdzenie}[Wyznacznik jako działanie na najwyższej potędze]
\label{thm:lin11-det-wedge}
Niech \(\dim V=n<\infty\) i
\(A\colon V\to V\) będzie
endomorfizmem. Przestrzeń
\(\Lambda^n(V)\) jest jednowymiarowa,
więc istnieje dokładnie jeden
skalar \(d\in K\) taki, że
\[
 \Lambda^n A=d\,
   \operatorname{id}_{\Lambda^n(V)}.
\]
Skalar ten jest wyznacznikiem
\(A\) określonym w
Rozdziale~\ref{ch:lin6-wyznaczniki}.
W szczególności
\(\Lambda^n A=(\det A)\operatorname{id}\),
również dla \(n=0\), kiedy
\(\det A=1_K\).
\end{twierdzenie}
\begin{proof}
Dla \(n\ge1\) wybierzmy
bazę \(B=(e_1,\ldots,e_n)\).
Z
Twierdzenia~\ref{thm:lin11-zew-baza}
wektor
\(\varepsilon=e_1\wedge\cdots\wedge e_n\)
stanowi bazę jednowymiarowej
przestrzeni \(\Lambda^n(V)\).
Istnieje więc dokładnie
jeden skalar \(d\) o
\((\Lambda^n A)\varepsilon=d\varepsilon\).
Jeśli macierz \(A\) w bazie \(B\)
ma współczynniki \(a_{ij}\),
to
Twierdzenie~\ref{thm:lin11-minory}
przy \(k=n\) mówi, że
\[
 d=\det(a_{ij})_{i,j=1}^n.
\]
Jest to wyznacznik
zdefiniowany wcześniej,
niezależny od bazy;
ponieważ \(d\) jest skalarem
działania na przestrzeni
jednowymiarowej, jego
niezależność od bazy
widać także bezpośrednio.
Dla \(n=0\) jedynym
endomorfizmem \(V=\{0\}\)
jest identyczność,
a \(\Lambda^0 A=\operatorname{id}_K\).
Wzór jest wtedy zgodny
z umową \(\det A=1_K\).
\end{proof}

\begin{figure}[!htbp]
\centering
\begin{tikzpicture}[>=Latex,scale=0.87]
 \coordinate (O) at (0,0);
 \coordinate (U) at (3,1);
 \coordinate (V) at (1,2);
 \coordinate (UV) at (4,3);
 \fill[manifoldblue!15] (O)--(U)--(UV)--(V)--cycle;
 \draw[manifoldblue,thick] (U)--(UV)--(V);
 \draw[->,gray] (-0.35,0)--(4.9,0) node[right] {\(e_1\)};
 \draw[->,gray] (0,-0.35)--(0,3.8) node[above] {\(e_2\)};
 \draw[->,very thick,manifoldred] (O)--(U)
   node[pos=.68,below right] {\(u=(3,1)\)};
 \draw[->,very thick,manifoldgreen] (O)--(V)
   node[pos=.68,left] {\(v=(1,2)\)};
 \fill[manifoldblue] (UV) circle (1.5pt)
   node[above right] {\(u+v\)};
 \node[fill=white,inner sep=2pt] at (2.45,1.78)
   {\(\det[u\ v]=5\)};
\end{tikzpicture}
\caption{W \(\mathbb R^2\) iloczyn zewnętrzny
\(u\wedge v=5e_1\wedge e_2\)
koduje zorientowane pole równoległoboku.
Zmiana kolejności wektorów zmienia znak.}
\label{fig:lin11-wedge-pole}
\end{figure}

\section{Potęgi zewnętrzne dualu i formy naprzemienne}
\label{sec:lin11-dual-formy}

\begin{twierdzenie}[Formy naprzemienne jako elementy dualu]
\label{thm:lin11-formy-alt}
Niech \(V\) będzie przestrzenią
wektorową nad \(K\), a \(k\ge0\).
Istnieje kanoniczny izomorfizm
\[
 (\Lambda^k(V))^*
 \cong \operatorname{Alt}^k(V^*),
\]
gdzie prawa strona oznacza
przestrzeń naprzemiennych form
\(k\)-liniowych \(V^k\to K\)
z Rozdziału~\ref{ch:lin8-wieloliniowe}.
Jeśli ponadto \(\dim V<\infty\),
istnieje kanoniczny izomorfizm
\[
 \mathcal E_k\colon
 \Lambda^k(V^*)\longrightarrow
 \operatorname{Alt}^k(V^*)
\]
zadany dla \(k\ge1\) na
wielowektorach prostych wzorem
\[
 \mathcal E_k(\varphi_1\wedge\cdots\wedge\varphi_k)
 (v_1,\ldots,v_k)
 =\det(\varphi_a(v_b))_{a,b=1}^k.
\]
Dla \(k=0\) przyjmujemy
\(\mathcal E_0=\operatorname{id}_K\).
Oba izomorfizmy obowiązują
nad ciałem dowolnej charakterystyki.
\end{twierdzenie}
\begin{proof}
Dla \(k=0\) lewa strona pierwszego
izomorfizmu jest równa \(K^*\), a prawa
\(K\). W tym przypadku izomorfizm
jest dany przez ewaluację
\(L\mapsto L(1_K)\), której
odwrotnością jest
\(a\mapsto[x\mapsto ax]\).
Dla \(k\ge1\) funkcjonał
\(L\in(\Lambda^k(V))^*\)
wyznacza odwzorowanie
\[
 F_L(v_1,\ldots,v_k)
 :=L(v_1\wedge\cdots\wedge v_k),
\]
które jest \(k\)-liniowe
i naprzemienne.
Odwrotnie, dla każdej
naprzemiennej formy \(F\)
Twierdzenie~\ref{thm:lin11-potega-wedge-uniwersalna}
daje jednoznaczny funkcjonał
\(L_F\) spełniający
\(L_F(v_1\wedge\cdots\wedge v_k)
=F(v_1,\ldots,v_k)\).
Oba przyporządkowania są
liniowe i wzajemnie odwrotne,
ponieważ wielowektory proste
rozpinają \(\Lambda^k(V)\).

Załóżmy teraz \(\dim V=n<\infty\).
Dla \(\varphi_1,\ldots,\varphi_k\in V^*\)
wyznacznik z tezy jest
liniowy oddzielnie w każdym
\(\varphi_a\) i znika,
gdy dwa z nich są równe,
bo wtedy dwa wiersze są
identyczne.
Wobec tego własność uniwersalna
potęgi zewnętrznej \(V^*\)
daje dobrze określone liniowe
odwzorowanie \(\mathcal E_k\).
Wartości są naprzemiennymi
formami na \(V\), ponieważ
wyznacznik jest także
naprzemienny i wieloliniowy
w kolumnach.

Wybierzmy bazę
\(B=(e_1,\ldots,e_n)\)
i jej bazę dualną
\((e^1,\ldots,e^n)\).
Dla rosnących krotek \(I,J\)
długości \(k\) mamy
\[
 \mathcal E_k(
 e^{i_1}\wedge\cdots\wedge e^{i_k})
 (e_{j_1},\ldots,e_{j_k})
 =\det(\delta_{i_a j_b})
 =\delta_{I,J}.
\]
Oznacza to, że baza
\(\Lambda^k(V^*)\) z
Twierdzenia~\ref{thm:lin11-zew-baza}
przechodzi na bazę form
naprzemiennych opisaną w
Twierdzeniu~\ref{thm:lin8-wymiar-alt}.
Odwzorowanie jest zatem
izomorfizmem.
Dla \(k>n\) obie przestrzenie
są zerowe, a dla \(k=0\)
izomorfizm jest identycznością.
\end{proof}

\begin{wniosek}[Iloczyn zewnętrzny form — wzór z przetasowaniami]
\label{cor:lin11-shuffle}
Niech \(V\) będzie skończenie
wymiarową przestrzenią nad \(K\).
Niech \(p,q\in\mathbb N_0\),
\(\alpha\in\operatorname{Alt}^p(V^*)\),
\(\beta\in\operatorname{Alt}^q(V^*)\).
Przenosząc mnożenie z
\(\Lambda(V^*)\) przez izomorfizmy
\(\mathcal E_p,\mathcal E_q,
\mathcal E_{p+q}\), otrzymujemy
iloczyn zewnętrzny form
o wzorze
\[
 \begin{aligned}
 &(\alpha\wedge\beta)(v_1,\ldots,v_{p+q})\\
 &\quad=\sum_{\sigma\in\operatorname{Sh}(p,q)}
  \operatorname{sgn}(\sigma)
  \alpha(v_{\sigma(1)},\ldots,v_{\sigma(p)})
  \beta(v_{\sigma(p+1)},\ldots,v_{\sigma(p+q)}),
 \end{aligned}
\]
gdzie \(\operatorname{Sh}(p,q)\) jest
zbiorem permutacji \(\sigma\in S_{p+q}\)
spełniających
\(\sigma(1)<\cdots<\sigma(p)\) oraz
\(\sigma(p+1)<\cdots<\sigma(p+q)\).
Dla \(p=0\) lub \(q=0\)
prawa strona jest zwykłym
mnożeniem przez skalar.
Formuła nie zawiera dzielenia
przez silnie i obowiązuje
w dowolnej charakterystyce.
\end{wniosek}
\begin{proof}
Wystarczy udowodnić wzór
dla \(\alpha\) i \(\beta\)
pochodzących od prostych
iloczynów elementów \(V^*\):
takie elementy rozpinają
\(\Lambda^p(V^*)\) i
\(\Lambda^q(V^*)\), a oba
wyrażenia są dwuliniowe.
Niech
\(\alpha=\mathcal E_p(
 \varphi_1\wedge\cdots\wedge\varphi_p)\),
\(\beta=\mathcal E_q(
 \psi_1\wedge\cdots\wedge\psi_q)\).
Wtedy \(\alpha\wedge\beta\)
jest obrazem
\(\varphi_1\wedge\cdots\wedge\varphi_p
\wedge\psi_1\wedge\cdots\wedge\psi_q\),
więc jego wartość na
\((v_1,\ldots,v_{p+q})\)
jest wyznacznikiem macierzy
o pierwszych \(p\) wierszach
\(\varphi_a(v_b)\) i ostatnich
\(q\) wierszach \(\psi_c(v_b)\).
Rozwińmy ten wyznacznik
według pierwszych \(p\) wierszy.
Wybór \(p\) kolumn
indeksowanych rosnąco,
oraz rosnąco indeksowanych
pozostałych \(q\) kolumn,
jest równoważny wyborowi
przetasowania \(\sigma\).
Znak rozwinięcia jest równy
\(\operatorname{sgn}(\sigma)\).
Odpowiednie dwa minory
są dokładnie wartościami
\(\alpha\) i \(\beta\)
na wskazanych argumentach.
Otrzymujemy wzór z tezy
dla form prostych, a
dwuliniowość daje go dla
dowolnych \(\alpha,\beta\).
Rachunek wyznacznika i
rozwinięcie Laplace'a są
poprawne nad dowolnym ciałem.
\end{proof}

\begin{przyklad}[Kontrola współczynnika w iloczynie form]
\label{prz:lin11-shuffle}
Niech \(V=\mathbb R^3\)
z bazą \((e_1,e_2,e_3)\)
i dualną bazą \((e^1,e^2,e^3)\).
Weźmy
\(\alpha=\mathcal E_2(e^1\wedge e^2)\)
i \(\beta=\mathcal E_1(e^3)\).
Na krotce \((e_1,e_2,e_3)\)
z trzech przetasowań \((2,1)\)
jedynie identyczność wnosi
niezerowy składnik, równy \(1\).
Stąd
\[
 (\alpha\wedge\beta)(e_1,e_2,e_3)=1.
\]
Jest to także
\(\det I_3\), zgodnie z
\(\mathcal E_3(e^1\wedge e^2\wedge e^3)\).
W odróżnieniu od
nienormalizowanej sumy po
wszystkich \(3!\) permutacjach
nie pojawia się czynnik \(2!\).
\end{przyklad}

\section{Symetryzacja i alternowanie tensorów}
\label{sec:lin11-projekcje}

\begin{stwierdzenie}[Permutacje czynników tensorowych]
\label{prop:lin11-permutacje}
Niech \(k\ge1\).
Dla każdej \(\sigma\in S_k\)
istnieje dokładnie jedno liniowe
odwzorowanie
\(P_\sigma\colon V^{\otimes k}\to V^{\otimes k}\)
takie, że
\[
 P_\sigma(v_1\otimes\cdots\otimes v_k)
 =v_{\sigma^{-1}(1)}\otimes\cdots
  \otimes v_{\sigma^{-1}(k)}.
\]
Spełnia ono
\(P_\sigma P_\tau=P_{\sigma\tau}\).
Operatory
\[
 \mathcal S_k=\sum_{\sigma\in S_k}P_\sigma,
 \qquad
 \mathcal A_k=\sum_{\sigma\in S_k}
                \operatorname{sgn}(\sigma)P_\sigma
\]
są liniowe nad każdym ciałem.
Zachodzą
\[
 \mathcal S_k^2=k!\,\mathcal S_k,\qquad
 \mathcal A_k^2=k!\,\mathcal A_k.
\]
Jeżeli \(k!1_K\ne0_K\),
to \((k!)^{-1}\mathcal S_k\)
i \((k!)^{-1}\mathcal A_k\)
są rzutowaniami liniowymi.
\end{stwierdzenie}
\begin{proof}
Zadany wzór jest \(k\)-liniowy
w \(v_1,\ldots,v_k\).
Istnienie i jednoznaczność
\(P_\sigma\) wynikają z
Twierdzenia~\ref{thm:lin10-wielokrotny}.
Dla tensora prostego po
zastosowaniu najpierw \(P_\tau\),
a potem \(P_\sigma\),
na miejscu \(i\) stoi
\(v_{\tau^{-1}(\sigma^{-1}(i))}
=v_{(\sigma\tau)^{-1}(i)}\).
Otrzymujemy
\(P_\sigma P_\tau=P_{\sigma\tau}\)
na tensorach prostych, a więc
na całej przestrzeni.

W podwójnej sumie
\(\mathcal S_k^2\)
dla każdej permutacji
\(\rho\in S_k\) występuje
dokładnie \(k!\) par
\((\sigma,\tau)\) o \(\sigma\tau=\rho\).
Stąd \(\mathcal S_k^2=k!\mathcal S_k\).
Dla \(\mathcal A_k\) dodatkowy
współczynnik
\(\operatorname{sgn}(\sigma)
\operatorname{sgn}(\tau)
=\operatorname{sgn}(\rho)\)
daje ten sam wynik.
Gdy \(k!1_K\ne0_K\),
jest on odwracalny w ciele.
Podzielenie obu równości
przez \((k!)^2\) dowodzi
idempotencji znormalizowanych
operatorów.
\end{proof}

\begin{uwaga}[Dlaczego iloraz jest konstrukcją podstawową]
\label{rem:lin11-kchar}
Znormalizowana symetryzacja
i alternowanie stanowią
wygodne narzędzia, gdy
charakterystyka ciała nie dzieli \(k!\).
Jeśli \(k!1_K=0_K\),
wzory z dzieleniem przez \(k!\)
nie mają sensu.
Nie oznacza to jednak, że
\(S^k(V)\) lub \(\Lambda^k(V)\)
przestają istnieć:
obie przestrzenie są
zdefiniowane jako ilorazy
algebry tensorowej, bez
jakiegokolwiek dzielenia.
Rozdział~\ref{ch:lin8-wieloliniowe}
opisuje tę samą przeszkodę
dla operatorów działających
na formach wieloliniowych.
\end{uwaga}

\begin{stwierdzenie}[Symetryczny iloraz a symetryzator]
\label{prop:lin11-projekcje-ilorazy}
Niech \(k\ge1\) i załóżmy,
że \(k!1_K\ne0_K\).
Niech \(q_k^S\colon
V^{\otimes k}\to S^k(V)\)
i \(q_k^\Lambda\colon
V^{\otimes k}\to\Lambda^k(V)\)
będą naturalnymi odwzorowaniami
ilorazowymi.
Wówczas ograniczenia
\[
 \begin{aligned}
 q_k^S|_{\operatorname{im}((k!)^{-1}\mathcal S_k)}
 &\colon \operatorname{im}((k!)^{-1}\mathcal S_k)
       \longrightarrow S^k(V),\\
 q_k^\Lambda|_{\operatorname{im}((k!)^{-1}\mathcal A_k)}
 &\colon \operatorname{im}((k!)^{-1}\mathcal A_k)
       \longrightarrow \Lambda^k(V)
 \end{aligned}
\]
są izomorfizmami.
Są to dwa różne sposoby
opisania tych samych ilorazów,
dostępne przy wskazanym
założeniu o charakterystyce.
\end{stwierdzenie}
\begin{proof}
Oznaczmy znormalizowane operatory
przez \(P_S\) i \(P_\Lambda\).
Dla permutacji czynników
mamy
\(q_k^S P_\sigma=q_k^S\)
z przemienności w \(S(V)\)
oraz
\(q_k^\Lambda P_\sigma
=\operatorname{sgn}(\sigma)q_k^\Lambda\)
z reguły znaków.
Zatem po zsumowaniu
\[
 q_k^S P_S=q_k^S,\qquad
 q_k^\Lambda P_\Lambda=q_k^\Lambda.
\]
Te równości zapewniają
surjektywność obu ograniczeń:
każdy element ilorazu
ma przedstawiciela z obrazu
odpowiedniego rzutowania.

Wykażmy teraz, że
\(P_S\) zeruje jądro \(q_k^S\).
Ideał symetryczny w stopniu
\(k\) jest rozpięty przez
tensory postaci
\(x\otimes(u\otimes v-v\otimes u)\otimes y\),
z możliwymi pustymi czynnikami
\(x,y\). Dla każdej z takich
różnic dwie sumy po
wszystkich permutacjach
czynników są jednakowe,
więc ich różnica jest zerowa.
Z liniowości \(P_S\)
zeruje całe jądro.
Analogicznie
\(P_\Lambda\) zeruje jądro
\(q_k^\Lambda\):
dla tensora z dwoma
sąsiednimi równymi wektorami
składniki sumy alternującej
łączą się w pary,
które mają równe wartości
tensorowe i przeciwne znaki.
Nie wymaga to dzielenia
przez \(2\).
Każdy element jądra
jest kombinacją takich
tensorów z dwoma
równymi sąsiednimi czynnikami.

Jeśli \(x=P_S(t)\)
i \(q_k^S(x)=0\), to
\(P_S(x)=0\), ponieważ
\(x\in\ker q_k^S\).
Z drugiej strony
\(P_S(x)=P_S^2(t)=P_S(t)=x\).
Stąd \(x=0\).
Ten sam argument
z \(P_\Lambda\) wykazuje
iniektywność drugiego
ograniczenia. Oba są więc
izomorfizmami.
\end{proof}

\section{Podsumowanie}
\label{sec:lin11-podsumowanie}

Algebra tensorowa \(T(V)\)
łączy wszystkie potęgi
\(V^{\otimes k}\) i zachowuje
kolejność czynników.
Jej iloraz \(S(V)\) wymusza
przemienność i w skończonym
wymiarze jest izomorficzny
z algebrą wielomianów.
Iloraz \(\Lambda(V)\) wymusza
\(v\wedge v=0\) i daje
naprzemienne iloczyny zewnętrzne.
W skończonym wymiarze
potęgi zewnętrzne mają bazy
indeksowane rosnącymi krotkami
i wymiary \(\binom nk\).
Współczynniki
\(v_1\wedge\cdots\wedge v_k\)
są minorami macierzy kolumn,
a macierz \(\Lambda^k A\)
jest złożona z minorów
macierzy \(A\).
Przy \(k=\dim V\) działanie
na jednowymiarowej
\(\Lambda^k(V)\) jest mnożeniem
przez wyznacznik.
Własność uniwersalna
potęgi zewnętrznej
identyfikuje jej dual z
przestrzenią form naprzemiennych,
a iloczyn form otrzymuje
wzór z przetasowaniami,
poprawny nad każdym ciałem.
Symetryzatory i antysymetryzatory
są wygodnymi dodatkowymi
reprezentacjami ilorazów
tylko przy założeniu, że
\(k!\) jest odwracalne w \(K\).
